% SPDX-License-Identifier: CC-BY-SA-4.0
% Author: Matthieu Perrin
% Part: <Nom de la partie>
% Section: <Nom de la section>
% Sub-section: <Nom de la sous-section>  % (facultatif, laisser vide si non utilisé)
% Frame: <Titre de la slide>

\begingroup

\begin{frame}{Confection de l'emploi du temps}

  \on[text, top=-3mm]{
    \Probleme{$k$-coloring (pour un entier $k \ge 2$)}{
      Un graphe $G = \langle V, E \rangle$ non-orienté
    }{
      Existe-t-il une coloration $c : V \rightarrow \{1, ..., k\}$ telle que :
      \begin{itemize}
      \item $\forall \{u, v\} \in E, c(u) \neq c(v)$ ?
      \end{itemize}
    }
  }

  \onBlock[y=1mm]{Théorème (admis)}{
    \begin{itemize}
     \item \textsc{$k$-coloring} est \textsc{np-}complet pour $k\ge 3$.
     \item \textsc{$2$-coloring} est polynomial.
    \end{itemize}
  }
  
  \onExampleBlock[bottom=3mm]{Exemple -- Emploi du temps}{
    \begin{description}[Sommets :]
    \item[Sommets :]  un ensemble de cours
    \item[Arêtes :]   relation de conflit
      \begin{itemize}
      \item même enseignant, groupe, salle, ... 
      \end{itemize}
    \item[Couleurs :] horaires dans la semaine
    \end{description}
  }

  \on[x=35mm, y=-13mm]{
    \begin{tikzpicture}
      \node[circle, draw=structure, fill=structure!20] (11) at (090:07mm) {}; 
      \node[circle, draw=example, fill=example!20]     (12) at (162:07mm) {}; 
      \node[circle, draw=example, fill=example!20]     (13) at (234:07mm) {}; 
      \node[circle, draw=alert, fill=alert!20]         (14) at (306:07mm) {}; 
      \node[circle, draw=alert, fill=alert!20]         (15) at (018:07mm) {}; 
      \node[circle, draw=alert, fill=alert!20]         (21) at (090:15mm) {}; 
      \node[circle, draw=structure, fill=structure!20] (22) at (162:15mm) {}; 
      \node[circle, draw=alert, fill=alert!20]         (23) at (234:15mm) {}; 
      \node[circle, draw=example, fill=example!20]     (24) at (306:15mm) {}; 
      \node[circle, draw=structure, fill=structure!20] (25) at (018:15mm) {}; 

      \draw (11) -- (13) -- (15) -- (12) -- (14) -- (11);
      \draw (21) -- (22) -- (23) -- (24) -- (25) -- (21);
      \draw (11) -- (21);
      \draw (12) -- (22);
      \draw (13) -- (23);
      \draw (14) -- (24);
      \draw (15) -- (25);
    \end{tikzpicture}
  }

  \footnoteref{Even, Itai, Shamir. \emph{On the Complexity of Timetable and Multicommodity Flow Problems}. SICOMP (1976)}

\end{frame}

\endgroup
